\section{Problem Definition}
\label{sec:problem}
This section defines the attribute synthesis problem and the correctness
requirements for its solutions. Given a rewrite's LHS metadata, the goal is
to decide whether valid RHS attributes exist and construct them when
feasible. Application also respects the user-written conditions. We also
specify the mathematical fragment for which we establish correctness and completeness.

\subsection{Tensor Expressions and Rewrites}
Let $\mathcal X=\{X_1,\ldots,X_n\}$ be the input tensor names. Each $X_i$ has
symbolic shape $s_i$; a match supplies its concrete shape. Let $\mathcal O$
be the operations. An operation $o\in\mathcal O$ has $k$ inputs and an
attribute descriptor $a$. Tensor expressions have the form
\begin{equation}
 T ::= X_i \mid o_a(T_1,\ldots,T_k).
\end{equation}
A rewrite is $L\Rightarrow R$. Its LHS attributes are known, and synthesis
constructs the unknown RHS axis fields. We use the following notation throughout:
\begin{equation}
\begin{aligned}
 s&=(s_1,\ldots,s_n) &&\text{input shapes},\\
 \alpha&=\text{LHS attributes},\\
 m&=(s,\alpha) &&\text{match metadata},\\
 \beta&=\text{completed RHS attributes}.
\end{aligned}
\label{eq:metadata}
\end{equation}
Match metadata $m$ is concrete at a match and makes every LHS operation legal.
Scalar types are fixed in the exact-real fragment.
Writing $X=(X_1,\ldots,X_n)$ for the input values, the two computations are
$L_m(X)$ and $R_{m,\beta}(X)$. An attribute builder reads $m$, never the
entries of $X$.

\subsection{Valid Attribute Assignments}
\label{sec:valid}
We define what counts as a correct synthesis result.
Let $\beta$ denote the completed RHS attributes.
$\operatorname{Legal}_R$ includes the legality of every intermediate
operation, $\operatorname{Shape}$ is an ordered shape vector, and $X$ ranges
over all tensor values with the input shapes specified by $m$. We define
\begin{equation}
\begin{aligned}
 \operatorname{Valid}(m,\beta)\iff{}&
 \operatorname{Legal}_R(m,\beta)\\
 &\land\operatorname{Shape}_L(m)=\operatorname{Shape}_R(m,\beta)\\
 &\land\forall X,\; L_m(X)=R_{m,\beta}(X).
\end{aligned}
\label{eq:valid}
\end{equation}
Thus, valid attributes make the RHS legal and preserve both the ordered
output shape and the values for every compatible input. Dtype compatibility
is assumed; axis synthesis and validation do not check dtype contracts.

Synthesis fills unknown axis fields from metadata $m$ and preserves supplied
fields. Supplied expressions must be defined whenever the RHS attribute
builder runs.
The resulting $\beta$ must be correct for every compatible input value,
so the quantifier order is $\exists\beta\,\forall X$.

\subsection{Synthesis and Execution Contract}
\label{sec:find}
For fully automatic axis fields in Section~\ref{sec:formal-scope}'s fragment,
we seek two metadata functions: the \emph{feasibility check}
\code{check} and the attribute builder \code{build}. Given match metadata
$m$, $\operatorname{check}(m)$ decides whether a valid RHS completion
exists, and $\operatorname{build}(m)$ constructs one when the feasibility
check succeeds. These functions must satisfy:
\begin{align}
 \operatorname{check}(m)&\iff\exists\beta:\operatorname{Valid}(m,\beta),
 \label{eq:check-goal}\\
 \operatorname{check}(m)&\Rightarrow
 \operatorname{Valid}(m,\operatorname{build}(m)).
 \label{eq:build-goal}
\end{align}
Both functions depend only on match metadata and must satisfy these
requirements for every legal instance in the defined fragment, across
varying ranks and shapes.
For a legal match in this fragment, the \emph{applicability condition}
requires both the written condition $C(m)$ and the feasibility check
$\operatorname{check}(m)$ to hold. Here $C$ is the written \code{where}
condition, or true when absent.

\paperparagraph{Execution contract.}
TEPL compiles the written condition, feasibility check, and attribute
builder once per rule. A legal match evaluates the written condition first,
then the feasibility check, and runs the builder only if both succeed.
A false condition skips the match without evaluating the remaining steps.
Matching evaluates metadata programs without repeating synthesis or proof
search. This execution order also applies outside the complete synthesis
fragment.

Within the domains in Table~\ref{tab:guarantee-map}, an accepted rule
produces defined attributes and a legal, equivalent replacement on every
legal match that passes both conditions. This guarantee assumes the
operation semantics and execution assumptions in
Section~\ref{sec:implementation-scope}.
Section~\ref{sec:method} establishes compilation correspondence, and
Section~\ref{sec:candidate-validation} establishes conditional validation
soundness.

\subsection{Formal Scope and Guarantees}
\label{sec:formal-scope}
The completeness theorem covers \emph{axis-structural multilinear}
computations: operations connect, reorder, and sum whole axes, while values
are computed through products and sums of independent input entries.
Multilinear means linear in each input when the others are fixed.
Synthesis chooses independent RHS \code{axis} fields under the restrictions
in Table~\ref{tab:direct-scope}. Extent and coordinate fields are not inferred.

\begin{table}[!htbp]
\caption{Restrictions of the direct completeness theorem and boundary examples.}
\label{tab:direct-scope}
\small
\begin{tabularx}{\linewidth}{@{}p{0.18\linewidth}YY@{}}\toprule
Condition & Restriction & Boundary example\\\midrule
\textbf{D1. Whole-axis structure}
& Connect, reorder, and sum whole axes in the same direction. No new axes,
coordinate splitting or merging, or additional weights.
& A slice taking $[a,b,c,d]$ to $[b,c]$ selects only part of an axis.\\
\addlinespace
\textbf{D2. Unrestricted RHS axis choices}
& All RHS axis fields are unknown. No fixed choices, list lengths, or layout
constraints beyond operation legality; empty lists are allowed.
& A dot with a fixed batch axis followed by a sum may be valid. Contracting
that axis early would violate the fixed attribute.\\
\addlinespace
\textbf{D3. Output ordering}
& Freely reorder axes, preserve their input order, or concatenate groups in
the declared order.
& A declaration interleaving two groups as $(i_1,j_1,i_2,j_2)$ requires an
additional ordering rule. This is a boundary of the current rules.\\
\addlinespace
\textbf{D4. Independent axis unknowns}
& Different RHS fields and nodes do not share one unknown assignment.
& To match $A+\operatorname{transpose}(B)$ for independent $2\times2$ inputs,
two RHS transposes require identity for $A$ and swap for $B$. A shared
permutation cannot satisfy both.\\\bottomrule
\end{tabularx}
\end{table}

\paperparagraph{Operation declarations.}
Each operation declares which input dimensions are connected, which are summed,
and the output order. For example, contraction pairs axes of equal extent
from two inputs and multiplies entries at matching coordinates. Paired axes
retained in the output are \emph{batch} axes; those summed out are
\emph{contraction} axes. Unpaired axes remain \emph{free}. The output order
is batch, left-free, then right-free, preserving each input's free-axis
order. Batch and contraction positions are distinct and in range within
each operand, and the left and right lists pair axes in list order.

\paperparagraph{Input and rewrite assumptions.}
\begin{itemize}
\item Inputs have finite ranks, and every input dimension has extent at least
two; scalars are allowed.
\item Input entries are independent real numbers, and all calculations are
exact. Every LHS operation is legal.
\item Each \emph{read-once component} is a contraction/permutation/unary-summation
tree using each tensor in a nonempty input set exactly once. Corresponding
LHS and RHS components use the same input set.
\item With addition, expand the LHS into $L=\sum_t L_t$ using linearity and
bilinearity. The RHS must already be an explicit sum $R=\sum_t R_t$ of
matching components. Each pair $L_t,R_t$ uses the same input set $U_t$;
different terms have different input sets, which may overlap.
\end{itemize}
Captured LHS attributes retain their values when expressions are expanded;
a shared RHS unknown must not be split into independent choices.

\paperparagraph{Guarantees and their boundaries.}
Table~\ref{tab:guarantee-map} summarizes the results and their assumptions.
A \emph{sound schema} is a rewrite pattern with a separately proved
attribute construction (Appendix~\ref{app:axis-extensions}). Its conditions
and agreement of supplied fields with constructed fields must hold under
the applicability condition.

\begin{table}[!htbp]
\caption{Formal guarantees and their scope. Extents are dimension sizes;
all rows require legal source instances and the stated operation semantics.}
\label{tab:guarantee-map}
\small
\begin{tabularx}{\linewidth}{@{}p{0.14\linewidth}Y>{\raggedright\arraybackslash}p{0.12\linewidth}Y@{}}\toprule
Result & Guarantee & Extents & Other assumptions and reference\\\midrule
Synthesis & Finds valid RHS axis assignments whenever they exist.
& $\ge2$ & D1--D4, read-once components, and exact arithmetic
(Theorem~\ref{thm:complete-synthesis}).\\
\addlinespace
Compilation & Compiled check and builder implement the mathematical functions.
& $\ge2$ & Admitted plans in the synthesis fragment
(Theorem~\ref{thm:symbolic-compilation}).\\
\addlinespace
Validation & Safe check evaluation and valid completed attributes under the applicability condition; supplied fields are allowed.
& $\ge0$ & Supported validation rules, exact arithmetic, and execution assumptions
(Theorem~\ref{thm:validator-soundness}; Section~\ref{sec:implementation-scope}).\\
\addlinespace
Sound schemas & Constructed replacements preserve legality, values, and ordered shape.
& $\ge0$, subject to schema conditions & Stated operation contracts and conditions for each schema
(Lemma~\ref{lem:certified-extensions}).\\\bottomrule
\end{tabularx}
\end{table}

Mixed explicit/inferred axis fields and zero or singleton extents have
soundness guarantees within the validation and schema domains, but no
completeness guarantee. Outside the complete synthesis fragment, a
feasibility check may reject feasible matches. Validation is sound but
incomplete: it can remain unresolved for a correct candidate. Successful
compilation does not establish that a rule has an applicable match.
Broader implementation support does not extend these formal guarantees;
Section~\ref{sec:scope-limitations} describes practical limitations.
